% Avsnitt 2.8, ur lectures/L02/README.md.

\section{Sammanfattning}
\label{c2:sec:review}

\needspace{6\baselineskip}
\subsection*{Det här ska du kunna nu}
\begin{itemize}
  \item Härleda en minimerad ekvation ur ett Karnaughdiagram och göra en VHDL-arkitektur av den.
  \item Läsa och skriva portar och signaler av typen \code{std_logic_vector}.
  \item Skriva en \code{process} med korrekt känslighetslista och ett \code{case} som täcker varje
    indatavärde.
  \item Bygga en design av mer än en \code{entity} med positionell \code{port map}, och säga
    varför två instanser av samma modul lägger två kopior av dess grindar på FPGA:n.
  \item Verifiera en modul genom att köra dess testbänk, och läsa ett misslyckat \code{assert}.
\end{itemize}

\needspace{6\baselineskip}
\subsection*{Frågor att testa dig själv med}
\begin{enumerate}
  \item Kan du, givet en sanningstabell, härleda dess ekvation både direkt och via ett
    Karnaughdiagram? Vad skiljer de två resultaten?
  \item Varför kan ett \code{case} inte stå direkt i en arkitekturkropp?
  \item Vad skyddar \code{when others} dig mot, när väljaren ”bara kan” vara \code{0} eller
    \code{1}?
  \item I \code{display1: entity work.display}, vad namnger \code{display1}, och vad namnger
    \code{work.display}? Skulle en andra instans kunna återanvända samma etikett?
  \item \code{hex_display}s arkitektur innehåller ingen logik alls. Var kommer grindarna i den
    färdiga designen ifrån?
  \item Du instansierade \code{display} två gånger i stället för att skriva dess \code{case} två
    gånger. Får FPGA:n en eller två kopior av den logiken? Vad vann du alltså egentligen?
  \item Boken skriver alltid \code{port map} positionellt. Vad går sönder om en entitets portar
    deklareras i en annan ordning än den testbänken förväntar sig, och när märker du det?
\end{enumerate}

\needspace{6\baselineskip}
\subsection*{Blicken framåt}
\Chapref{c3:ch} går vidare med:
\begin{itemize}
  \item Sekvensnät: kretsar som minns.
  \item D-låset, och D-vippan byggd av ett par av dem.
  \item Klockning, och när en signaltilldelning faktiskt får effekt.
  \item Flankdetektering: en knapptryckning som växlar en lysdiod.
\end{itemize}
